\documentstyle[aps,prb%
,preprint%
%,twocolumn%
]{revtex}
\preprint{Submitted to {\it Phys. Rev. B}}
\setlength{\textheight}{25cm}  % comment out for galley proofs
%\let\labelm\label \renewcommand{\label}[1]{\fbox{\tt #1}\labelm{#1}}
\begin{document}
\draft
\title{DRY FRICTION IN THE FRENKEL-KONTOROVA-TOMLINSON MODEL:
   \\I. STATIC PROPERTIES}
\author{Michael Weiss and Franz-Josef Elmer}
\address{Institut f\"ur Physik, Universit\"at
   Basel, CH-4056 Basel, Switzerland}
\maketitle
\begin{abstract}                                                     
We investigate wearless friction in a simple mechanical model called
Frenkel-Kontorova-Tomlinson model. It combines the Frenkel-Kontorova
model (i.e., a harmonic chain in a spatially periodic potential) with
the Tomlinson model (i.e., independent oscillators connected to a
sliding surface in a fixed potential describing the other surface).
We investigate static properties like the ground state, the
metastable states, and static friction, as well as the the kinetic
friction in the limit of quasistatic sliding.  Like in the
Frenkel-Kontorova model the behavior strongly depends on whether the
ratio of lattice constants is commensurate or incommensurate.  In the
incommensurate case Aubry's transition by breaking of analyticity
also appears in the Frenkel-Kontorova-Tomlinson model. The behavior
depends strongly on the strength of the interaction between the
sliding surfaces.  For increasing interaction, we find three
thresholds which denote the appearance of static friction, of kinetic
friction in the quasistatic limit, and of metastable states in that
order. These are identical only in the incommensurate case. In the
commensurate case static friction can be nonzero even though the
kinetic friction vanishes for sliding velocity going to zero.
\end{abstract}
\pacs{PACS numbers: 46.30.Pa, 05.70.Ln, 81.40.Pq%
%\\Submitted to Phys. Rev. B.%
}

\narrowtext

\section{\protect\label{INT}Introduction}

Dry friction is a phenomenon of everyday life. Since Coulomb's work
its basic phenomenological laws are well-known\cite{bow.54}: (i) The
friction force is independent of the area of the sliding surface.
(ii) It is proportional to the load. (iii) The kinetic friction,
i.e., the force to keep a body sliding at a constant velocity, does
not depend on the velocity and it is less than or equal to the static
friction, i.e., the force to start sliding. On a macroscopic level,
these laws are well understood in terms of the Bowden-Tabor adhesion
model\cite{bow.54,sin.92} which is a macroscopic model based on the
elastic and plastic properties of the sliding bodies.

In spite of the simplicity of Coulomb's laws, the sliding of two solid
bodies is a very complex phenomenon which operates mostly far away
from thermal equilibrium. It involves processes on various spatial
and temporal scales, from microscopic to macroscopic (for an overview
of the state of the art see Ref.~\onlinecite{sin.92}). Also
deviations from Coulomb laws have been found, which depend on the
material of the sliding bodies, the surface properties, the sliding
history, and the mechanical environment. Up to now no generally
accepted theory exists which is able (i) to explain these deviations
and (ii) to calculate friction forces from the bulk and surface
properties of the sliding bodies. However, in recent years modern
experimental technologies have made it possible to study wearless
friction between clean and atomically flat surfaces\cite{sin.92}.
There is some hope that theoretical models will lead to an
understanding of such less complex systems.

The first attempt to explain Coulomb's laws on the atomic level was
given by Tomlinson\cite{tom.29} in his pioneering work in 1929. He
considered the surface atoms as single independent oscillators which
are ``plucked'' by the atoms of the other surface like a guitar
string.  Even quasistatic sliding (i.e., sliding at an
infinitesimally small velocity) will lead to plucking of atoms if the
stiffness of the interaction between a surface atom and the bulk is
smaller than the stiffness of the interaction between this atom and
an atom from the other surface at the moment of closest
contact\cite{mcc.89,mcc.92}. During plucking the atom is assumed to
jump abruptly from one equilibrium position to another. Such jumps
lead to vibrations of the jumping atom. The kinetic energy of the
vibrating atom is assumed to dissipate totally into the bulk of the
sliding bodies due to excitation of some kind of waves. Thus finite
friction is possible even in the limit of zero sliding velocity,
contrary to viscous friction which vanishes in the quasistatic limit.
The assumptions that the jumps occur instantaneously and that the
atoms are uncoupled, i.e., a vibrating atom does not excite
vibrations of other atoms, lead to Coulomb's third
law\cite{tom.29,mcc.92}.

The easiest model for taking into account the coupling between atoms
is the Frenkel-Kontorova (FK) model\cite{fre.38} which is a model of
an adsorbed monolayer on an atomically flat surface. It is a
one-dimensional model with a chain of adsorbate atoms coupled
linearly by nearest-neighbor interactions. The chain interacts with a
spatially periodic potential. The FK model has also been used as a
simple friction model\cite{sok.84,mcc.89,mcc.92}.  The static
properties of this model strongly depend on the ratio of the lattice
constants of the adsorbate layer and the substrate surface. Aubry has
shown\cite{aub.78} that in the case of an irrational ratio, the
ground state can be shifted by an infinitesimally small force as long
as the strength of the periodic potential is below a critical value
(the point of analyticity breaking). Thus the static friction is
zero. This is true only in the thermodynamic limit, i.e., in the case
of an infinite number of adsorbate atoms. 

There is some ambiguity in the literature concerning the meaning of
the term {\em frictionless\/}. In the context of the FK model
frictionless means zero static friction. Another meaning of
frictionless found in the literature is a vanishing kinetic friction
in the limit of infinitesimally slow sliding\cite{mcc.89}. Recently
this notion has been called {\em
superlubricity\/}\cite{shi.93,hir.93}, which is a misleading term
because it does {\em not\/} imply that a {\em finite\/} sliding
velocity exists below which the kinetic friction is zero. It should
be emphasized that the two definitions of frictionless are {\em
not\/} equivalent. Of course, zero static friction implies zero
kinetic friction if we believe in Coulomb's third law. But a
vanishing kinetic friction does not imply a vanishing static
friction.  This can be easily seen in the Tomlinson model: When the
interaction of the oscillator with the sliding surface is weak, it
will not be plucked. Thus it can be moved adiabatically without
dissipation and the kinetic friction is zero. Nevertheless the
surface {\em pins\/} the oscillator. Therefore a finite force is
necessary to depin it which means that the static friction is not
zero. It is equal to the amplitude of the oscillating sliding force
in the case of adiabatic sliding.

The main disadvantage of the FK model is that the atoms are not
coupled to the sliding body. One simple way to overcome this
disadvantage is to couple each atom harmonically to a rigid body (see
Fig.~\ref{f1}).  The resulting model we call the
Frenkel-Kontorova-Tomlinson (FKT) model because it is a combination
of the FK model and the Tomlinson model.

The FKT model is similar to the Burridge-Knopoff (BK)
model\cite{bur.67}, which was proposed as a model of an earthquake
fault. The main difference is that in the BK model the interaction
with the lower body is replaced by a phenomenological dry friction
law. Usually a velocity-weakening law is chosen\cite{car.89}. For
small sliding velocities, this model exhibits avalanches (i.e., 
earthquakes). For that reason the BK model has become a popular model
for the investigation of self-organized
criticality\cite{car.89,sor.92}. Also soliton-like behavior has been
found\cite{sch.93}. 

This paper is the first part of a series of papers investigating the
FKT model. This part considers only the static properties. The
forthcoming parts will be devoted to kinetic friction and stick-slip
motion, respectively. This first part is organized as follows: 
Section~\ref{MOD} introduces the FKT model. Sec.~\ref{F0}
investigates the FKT model in the case of zero driving force. In
Sec.~\ref{STA} we calculate the static friction. The kinetic friction
in the quasistatic limit is considered in Sec.~\ref{KIN}. We conclude
with Sec.~\ref{CON}.


\section{\protect\label{MOD}The model}

The Frenkel-Kontorova-Tomlinson (FKT) model is a one-dimensional
lattice model for the atomic monolayer of the surface of a soft body
(upper body in Fig.~\ref{f1}) which slides on a hard body (lower body
in Fig.~\ref{f1}). The monolayer is described by a chain of $N$
particles with harmonic nearest-neighbor interactions (coil springs).
The interaction of each particle with the otherwise rigid upper body
is also harmonic (leaf springs). The equilibrium positions of the
particles due to this interaction define a regular lattice where the
lattice constant is assumed to be the bulk lattice constant of the
upper body. The interaction of the particles with the lower body is
described by a spatially periodic external potential which defines a
hard surface. The lower body is assumed to be fixed whereas the upper
body is movable. The model assumes motions only {\em parallel\/} to
the sliding surface.

The potential energy of the FKT model is \begin{eqnarray}
  \lefteqn{V(\xi_1,\cdots,\xi_N,x_B)=\frac{1}{2}
    \sum_{j=1}^N(\xi_j-\xi_{j-1})^2}\nonumber\\
   &&+\frac{\kappa}{2}\sum_{j=1}^N\xi_j^2
  +\frac{b}{2\pi}\sum_{j=1}^N\cos 2\pi (x_B+cj+\xi_j)-Fx_B.
  \label{MOD.V} 
\end{eqnarray} 
Here $c$ is the lattice constant of
the upper body, $x_B$ is the position of the upper body relative to
the lower surface, $\xi_j$ is the position of particle $j$ relative
to the support $x_B+cj$ of its leaf spring, $\kappa$ is the stiffness
of the leaf springs, $b$ is the strength of the external potential
which models the interaction with the lower body, and $F$ is the
force applied to the upper body.  All variables and parameters are
measured in dimensionless units.  They are based on the following
independent basic units: The length unit is the surface lattice
constant of the lower body, and the unit of the interaction strength
is the stiffness of the nearest-neighbor coupling. All other units
can be expressed in terms of these basic units. 

We choose periodic boundaries
\begin{equation}
  \xi_{j+N}=\xi_j.
  \label{MOD.pb}
\end{equation}
This implies a rational upper lattice constant $c$ because $cN$ has
to be an integer
\begin{equation}
  c=\frac{P}{Q},
  \label{MOD.cpb}
\end{equation}
where $P$ and $Q$ are coprime, and $Q$ is a divisor of $N$.

For $F\not=0$ the potential energy $V$ is not bounded from below.
Thus no ground state exists and all local minima correspond to
metastable states. Therefore, the system is always in a
non-equilibrium state.  At finite temperature, the system can escape
from a metastable state to another metastable state with lower
energy. Thus the system creeps due to thermal
activation\cite{hes.94}. In this paper we restrict ourselves to zero
temperature.

The FKT model is invariant against the following transformations:
\begin{mathletters}
\label{MOD.sym}
\begin{itemize}
\item {\bf Translations} by an integer multiple of the lattice
constant of the lower body generated by
\begin{equation}
  x_B\to x_B+1.
  \label{MOD.Tsym}
\end{equation}
\item {\bf Cyclic permutations} of the deformations $\xi_j$ generated
by
\begin{equation}\begin{array}{rcl}
   x_B&\to &x_B+c,\\
   \xi_j&\to &\xi_{j+1}.
   \end{array}\label{MOD.CPsym}
\end{equation}
\item {\bf Reflection} of the upper body at $x=0$ 
\begin{equation}\begin{array}{rcl}
   x_B&\to &-x_B,\\
   \xi_j&\to &-\xi_{-j},\\
   F&\to &-F.
   \end{array}\label{MOD.I2sym}
\end{equation}
\item {\bf Rescalings} of the lattice constant, generated by
\begin{equation}
  c\to c+1.
  \label{MOD.Ssym}
\end{equation}
\item {\bf Inversion} of the counting order
\begin{equation}\begin{array}{rcl}
   c&\to &-c,\\\xi_j&\to &\xi_{-j}.
   \end{array}\label{MOD.I1sym}
\end{equation}
\end{itemize}
\end{mathletters}
These transformations leave the energy landscape invariant except 
(\ref{MOD.Tsym}) and (\ref{MOD.CPsym}) shift the potential by a 
constant value proportional to $F$. The symmetries (\ref{MOD.sym}) 
hold also for an arbitrary periodic potential except the reflection
symmetry (\ref{MOD.I2sym}) which holds only for potentials which are
invariant against reflection like the cosine potential in
(\ref{MOD.V}). The symmetries permit us to restrict $c$ to the
interval $[1/2,1]$.

For $\kappa=0$ the FKT model becomes the FK model.  For
$\kappa,b\to\infty$ but $b/\kappa$ finite, or similarly by dropping
the nearest-neighbor interaction, the FKT model turns into the
Tomlinson model of independent oscillators.

Similar to the FK model we expect a strong dependence of the FKT
model on whether $c$ is rational or irrational. There are 
two extreme cases: (i) the most commensurate case, $c=1$, and (ii)
the most incommensurate case given by the golden mean
$c=(\sqrt{5}-1)/2=0.618\ldots$ In order to treat the latter case 
we have to investigate a series of FKT models
with successively increasing numbers of particles. Since continued
fractions are the optimum rational approximations of irrational
numbers, we take successive pairs from the Fibonacci series
$1,1,2,3,5,8,13,\ldots$ in order to approximate the golden mean.

\section{\protect\label{F0}The FKT model without driving}

In this section we investigate the FKT model where the
external force is absent, i.e., $F=0$. Stationary states are the
extrema of (\ref{MOD.V}) which are solutions of
\begin{mathletters}
\label{F0.eq}
\begin{eqnarray}
 0=-\partial_{\xi_j}V&=&\xi_{j+1}+\xi_{j-1}-(2+\kappa)\xi_j
   \nonumber\\&&+b\sin 2\pi(x_B+cj+\xi_j) \label{F0.eqj}\\
 0=-\partial_{x_B}V&=&\kappa\sum_{j=1}^N\xi_j\label{F0.xisum}.
\end{eqnarray}
\end{mathletters}
It is important to note that $x_B$ is not a free variable because a
state which is stationary with respect to the $\xi_j$ need not be
stationary with respect to $x_B$.

Because of the symmetries (\ref{MOD.Tsym}), (\ref{MOD.CPsym}), and
(\ref{MOD.I2sym}) any solution can be generated from a solution
with
\begin{equation}
 x_B\in \left[0,\frac{1}{2Q}\right].
  \label{F0.xB.int}
\end{equation}

For $F=0$, the potential energy (\ref{MOD.V}) of the FKT model is
bounded from below. Thus a {\em ground state\/} exists. It strongly 
determines the static friction, as we will see below.

\subsection{\protect\label{F01}The ground state}

In order to calculate the ground state, we start with the trivial
case $b=0$ where the energy surface is a paraboloid in the
$N$-dimensional configuration space defined by $\{\xi_j\}$.
Eq.~(\ref{F0.xisum}) defines a plane in this configuration space on
which the solution is situated. The absolute minimum of the parabolic
energy surface is at $\xi_j=0$. This solution automatically fulfills
(\ref{F0.xisum}).  Thus $x_B$ is arbitrary.

In the general case, $b\not=0$, the external potential adds some
corrugation onto the parabolic energy surface which shifts the
absolute minimum a certain amount outside the origin. By changing
$x_B$, this minimum can be moved into the plane defined by
(\ref{F0.xisum}). Thus, the ground state position $x_B^G$ of the
upper body is no longer arbitrary. The deformations $\xi_j^G$ of the
absolute minimum in any direction cannot exceed half the periodicity
of the external potential:
\begin{equation}
  |\xi_j^G|\le \frac{1}{2},\quad j=1,\ldots,N.
  \label{F0.GS1}
\end{equation}
For small values of $b$ the $\xi_j^G$'s are linear in $b$. For
$b\to\infty$ the particles are situated in the local minima of the
external potential. In order to minimize the potential energy in the
springs, each particle will presumably sit in that potential well
which is directly underneath the support of the leaf spring (see
Fig.~\ref{f1}). Thus we conjecture that for the ground state
\begin{equation}
  \mbox{Int}\,(x_B^G+cj+\xi_j^G)=\mbox{Int}\,(x_B^G+cj),\quad 
   j=1,\ldots,N
  \label{F0.GS2}
\end{equation}
holds. The same property holds for the ground state of the FK
model\cite{aub.78}, where $c$ corresponds to the winding number.
Numerically we have not found any example demonstrating the opposite.
Furthermore, we found that any other state violates (\ref{F0.GS2}).
Thus the ground state is uniquely defined by (\ref{F0.GS2}). Note
that (\ref{F0.GS1}) and (\ref{F0.GS2}) are independent conditions.
The ground state is either {\em commensurate\/} or {\em
incommensurate\/} depending on whether $c$ is rational or irrational.

Aubry has shown\cite{aub.78} that the ground state of the 
FK model can be uniquely described by a {\em hull function\/}. 
This concept can also be applied to the FKT model. Thus we write
\begin{equation}
  \xi_j^G=g_S(x_B^G+cj).
  \label{F0.gS}
\end{equation}
From Eq.~(\ref{F0.eqj}), the hull function $g_S$ is a solution of the
delay equation
\begin{eqnarray}
  0&=&g_S(x+c)+g_S(x-c)-(2+\kappa)g_S(x)\nonumber\\
   &&+b\sin 2\pi[x+g_S(x)]. \label{F0.gS.eq}
\end{eqnarray}
where $x = x_B^G+cj$. Assuming that the ground state does not break
the symmetries (\ref{MOD.sym}) we get 
\begin{mathletters}
\label{F0.gS.prop}
\begin{equation}
  g_S(x+1)=g_S(x)=-g_S(-x).
  \label{F0.gS.sym}
\end{equation}
The properties (\ref{F0.GS1}) and (\ref{F0.GS2}) are equivalent to
\begin{equation}
   |g_S(x)|\le 1/2
  \label{F0.gS.p1}
\end{equation}
and
\begin{equation}
   |g_S(x)+x\,\mbox{mod}\,1-1/2|\le 1/2,
  \label{F0.gS.p2}
\end{equation}
respectively.
\end{mathletters}

Which values of $x_B$ fulfill Eq.~(\ref{F0.xisum})?
First we consider the commensurate case. Because of the symmetries
(\ref{MOD.sym}), every $x_B$ with $x_B\,\mbox{mod}\,(1/2Q)=0$ yields
a stationary state. Numerically we find that the ground state is
given by
\begin{equation}
 x_B^G=\left(\frac{1}{2}+n\right)\frac{1}{Q},\quad n\ {\rm integer}.
   \label{F0.xB}
\end{equation}

In the incommensurate case $x_B$ is arbitrary.  We can interpret
$x_B$ as a {\em phase\/} which is discrete for rational values of $c$
and continuous for irrational ones. Below we will see that this
interpretation becomes very important for the construction of
elementary defect solutions.

Figure~\ref{f.hull} shows the hull functions of numerically
calculated ground states for an incommensurate value of $c$.  We
clearly see that they fulfill the properties (\ref{F0.gS.prop}).
Fig.~\ref{f.hull} also shows the phenomenon of {\em breaking of 
analyticity\/}, first discovered by Aubry\cite{aub.78} in the FK
model. It means that in the incommensurate case the hull function is
no longer analytic for $b$ above a certain critical value $b_c^S$.
The nonanalyticity is caused by jumps which appear at certain points.
The biggest jump is at $x=0$. It means that a kind of forbidden
region has emerged around the maxima of the external potential where
no particle is allowed to be. Due to the nearest-neighbor interaction
jumps also appear at $x=cn$, $n$ integer. The size of the jumps
decreases with increasing order $|n|$. Since the hull function is a
periodic function and $c$ is an irrational number, jumps occur in any
neighborhood of any point $x\in[0,1]$.

In the Tomlinson limit of the FKT model (i.e., $\kappa,b\to\infty$)
we can discuss the hull function analytically
because we can drop the delay terms of (\ref{F0.gS.eq}). We get
the hull function $g_S(x)$ in parametric form
\begin{equation}
g_S(\tau)=\frac{b}{\kappa}\sin 2\pi\tau,\quad x(\tau)=\tau-g_S(\tau).
  \label{F0.g}
\end{equation}
It is uniquely defined as long as $b<\kappa/(2\pi)$. For
$b>\kappa/(2\pi)$ three different values of $g_S$ are possible at
least around $x=0$. The additional two values corresponds to a
metastable state and an unstable saddle of a single Tomlinson
oscillator (see Fig.~\ref{f.hull2}). Thus they are not relevant for
the hull function which describes the ground state. Furthermore, they
do not fulfill (\ref{F0.gS.p2}).

This analysis suggests that the main jump and all additional jumps of
the hull function (which are infinitesimally small in the Tomlinson
limit) are caused by the appearance of {\em metastable states\/}. In
the following subsection we therefore discuss the question what kind
of metastable states appear first and for which values of $b$?

\subsection{\protect\label{F02}Metastable states}

For small values of $b$ the corrugation is weak and the ground state
is still the only stationary stable state. If $b$ is larger than a 
critical value $b_c^K$, additional stationary states emerge which
have the following properties. They violate condition (\ref{F0.GS2}).
The matrix $\partial^2V/\partial\xi_i\partial\xi_j$ is positive
definite which means that these states are stable against
infinitesimal small disturbances in $\xi_j$, $j=1,\ldots,N$. But they
may be still unstable against infinitesimal small perturbations
which includes also disturbances in $x_B$. If $b$ exceeds a further
threshold $b_c^m\ge b_c^K$, the first of these states become stable. 

The first metastable state will not be very different from the ground
state. It can be described by the ground state plus a {\em defect\/}.
But what kind of defect?  Naively one might expect a defect caused by
a single hopping of one particle to its neighboring potential well.
For the FK model it is known that this is not the case\cite{val.88}.
The elementary defect in the FK model and in the FKT model too --- as
we will see below --- is a {\em phase kink\/} which separates two
domains with ground states corresponding to different phases $x_B$
defined by (\ref{F0.xB}) (see Fig.~\ref{f.kink}). The phase of each
domain is a kind of ``virtual'' body position which would be the real
position if the other domain does not exist.  In the FK model such a
phase kink is a localized deviation from the ground state which means
that the kink energy does not depend on the size $N$ of the system if
$N$ is much larger than the size of the kink.  On the other hand, in
the FKT model the energy of such a defect scales with $N$ because of
the {\em nonlocal\/} coupling of each particle via the upper body.
The nonlocality is expressed by condition (\ref{F0.xisum}) which
means that the real position $x_B$ is the average over the virtual
ones weighted by the domain sizes (see Fig.~\ref{f.kink}). Thus the
difference between the real $x_B$ and the virtual $x_B$ shifts {\em
each\/} particle out of its ground-state position even if it is far
away from the kink.

Another consequence of the nonlocal coupling is that we cannot treat
an infinitely extended FKT model because the ratio of the domain
sizes would not be defined. Thus we consider a very large but finite
system. Because of the periodic boundaries, the numbers of kinks and 
antikinks have to be equal.

The above-mentioned defect of one particle hopping into the
neighboring potential well corresponds to a pair of a kink and an
antikink which are together as closely as possible. In the following
we argue that the first metastable state consists of a pair of a kink
and an antikink which are as far away as possible, not because of the
usual exponentially decreasing kink-antikink interaction but because
of the nonlocality. The difference between the virtual values of
$x_B$ of the domain states is $1/Q$.

First we treat the Tomlinson limit where the kink-antikink
interaction disappears because of the absence of the nearest-neighbor
interaction. For $b>\kappa/(2\pi)$ and $x_B$ fixed, each Tomlinson
oscillator develops additional stationary states because the hull
function $g_S(x)$ defined by (\ref{F0.g}) looses its uniqueness. One
of these additional states is metastable for $x_B$ fixed (see
Fig.~\ref{f.kink}). Nevertheless for $b$ just above $\kappa/(2\pi)$
any state other than the ground state which fulfills the nonlocality
condition (\ref{F0.xisum}) is still unstable when $x_B$ is free.  In
order to see this the most commensurate case (i.e., $c=1$) is treated
in Appendix~\ref{AP1}. It turns out that the first metastable
kink-antikink state (see Fig.~\ref{f.kink}) is indeed that one with
equally sized domains (see Fig.~\ref{f.c1kink}). It appears for
$b>\kappa/4$. Thus 
\begin{equation}
  b_c^m=\frac{\kappa}{4},\quad{\rm for}\quad c=1.
  \label{F0.bcm.c0}
\end{equation}
This threshold also holds in general, not only in the Tomlinson
limit, because the nearest-neighbor interaction only leads to a finite
kink size, but does not change the domain states.

For other values of $c$ we have found numerically that the state with
equally sized domains is always the first metastable state. It has
the following properties:  (i) $x_B\,\mbox{mod}\,(1/Q)=0$. (ii) All
$\xi_j$'s fulfill (\ref{F0.GS1}).  (iii) All $\xi_j$'s fulfill
(\ref{F0.GS2}) except for those for which the support of their leaf
spring (see Fig.~\ref{f1}) is directly above a maximum of the
external potential. The displacement $\xi_j$ is positive in one
domain and negative in the other one.  Thus in one domain every
$Q$-th particle violates (\ref{F0.GS2}).

In the incommensurate case metastable states appear beyond the point
of breaking of analyticity, i.e., $b_c^m\ge b_c^S$.  Numerically we
have some evidence that $b_c^m=b_c^S$ always holds, as in the FK
model\cite{aub.78}.

\section{\protect\label{STA}Static friction}

Before investigating the static friction we have to
clarify the meaning of static friction in deterministic
zero-temperature models, like the FKT model. We do this from the 
point of view of nonlinear dynamics. 

\subsection{\protect\label{STA1}Definition of static friction}

From the phenomenological point of view the static friction $F_S$ is
defined as the smallest driving force $F$ which initiates sliding.
That is, any force $F$ below $F_S$ does not lead to a relative motion
of the surfaces. Thus, the static friction is defined by the {\em
boundaries\/} of $F$ for which the stationary, motionless {\em
state\/} is {\em stable\/}. Therefore the static friction depends on
the internal state of the real contact area between the sliding
surfaces which in the FKT model is determined by the configuration of
the $N$ particles describing the surface atoms of the upper body (see
Fig.~\ref{f1}). Since different states may have different stability
boundaries the static friction is in general not uniquely defined. It
will depend on the {\em history\/} of the friction system, a
well-known experimental fact. For example, it was
found\cite{hes.94,rab.65} that the static friction increases with the
sticking duration, i.e., the time during which the sliding surfaces
do not move (at least not on a macroscopic level).

In view of the non-uniqueness of the static friction the question
arises:  Is a sound definition of static friction possible which does
not depend on the history of the total friction system? The answer of
this question is positive and the appropriate definition reads: {\em
The static friction $F_S$ is the smallest driving force above which
no stable stationary state exists\/}. This definition is very similar
to the definition of the depinning force in models describing the
pinning-depinning transition of charge-density waves or interacting
vortices in typ-II superconductors.  In these models the dynamics is
usually treated in the overdamped limit which has the following
consequences. Assuming we start with a metastable state and
slowly increase $F$. The state follows $F$ adiabatically until a
certain value is reached where it annihilates with an unstable state
in a saddle-node bifurcation. Thus the system depins locally and some
degrees of freedom start to move. In the case of overdamped dynamics
the system will be trapped by the next relative minimum of the energy
surface in configuration space. Thus for a slowly increasing
driving force $F$ the system will follow a sequence of depinning
events where it locally relaxes into another pinned state. 
Eventually the system globally depins. This happens at least when the
last state disappears.

Applied to friction we see that the above definition of the static
friction makes sense in the absence of inertia effects both in the
internal degrees of freedom and in the macroscopic sliding body. It
is also independent of the history of the friction system. The local
reconfigurations may be interpreted as {\em
microslips\/}\cite{rab.65} which do not lead to sliding on a
macroscopic scale. It must, however, be realized that the study of a
massless sliding body is somewhat academic, and a discussion of
inertia effects will be crucial for a realistic treatment of static
friction. Here we emphasize the main difference between models for
pinning-depinning transitions, like the FK model, and models like the
FKT model, namely the {\em existence of a macroscopic sliding mass
which interacts directly with many or even all internal degrees of
freedom\/}, which is also responsible for the {\em nonlocal\/}
coupling of the internal degrees of freedom. For finite inertia we
expect that an unstable state does not evolve into another stable
state but rather leads to global sliding. Nevertheless, $F_S$ gives
an upper bound for the actual static friction. It is calculated in the
next subsection. The static friction in the more realistic case of
finite inertia will be discussed in the context of stick-slip motion
in part III.

\subsection{\protect\label{STA2}Static friction of the FKT model}

In the previous subsection we have defined the maximum static
friction $F_S$ as the driving force $F$ where the last stationary
state disappears. The stationary states are the solutions of
\begin{mathletters}
\label{STA.eq}
\begin{eqnarray}
 0=-\partial_{\xi_j}V&=&\xi_{j+1}+\xi_{j-1}-(2+\kappa)\xi_j
   \nonumber\\&&+b\sin 2\pi(x_B+cj+\xi_j) \label{STA.eqj}\\
 0=-\partial_{x_B}V&=&\kappa\sum_{j=1}^N\xi_j+F\label{STA.eqB}.
\end{eqnarray}
\end{mathletters}

Without solving (\ref{STA.eq}) it is possible to give an overall upper
bound $F_S^{max}$ for the static friction $F_S$.  Taking the sum over
(\ref{STA.eqj}) and using (\ref{STA.eqB}) we get 
\begin{equation}
  F=-b\sum_{j=1}^N\sin 2\pi(x_B+cj+\xi_j),
  \label{STA.F}
\end{equation}
from which immediately follows that $|F|$ is always smaller than
\begin{equation}
  F_S^{max}=Nb.
  \label{STA.FSmax}
\end{equation}

Starting with a stable, stationary state of the FKT model without
driving, what happens if $F$ is quasistatically switched on? First,
$x_j$ and $x_B$ will change adiabatically, i.e., they are smooth
functions of $F$.  Eventually $F$ will reach a value -- the depinning
force -- where the state annihilates with an unstable stationary
state in a saddle-node bifurcation. From physical intuition and from
our numerical studies we are convinced that {\em the ground state of
the FKT model has the largest depinning force which therefore defines
the maximal static friction force $F_S$\/}. We are able to prove this
conjecture only for $c=1$ (see App.~\ref{AP0}), but it is consistent
with the above-mentioned experimental observations that the static
friction always increases with the sticking time. At the beginning of
sticking just after a global sliding event the contact area will be
in a metastable state.  On a large time scale where thermal agitation
becomes important, the system will overcome the barrier and move into
a more stable state. By repetition of this process it will eventually
reach the ground state.

For $c=1$ and $c=1/2$ we are able to calculate $F_S$ analytically (see
App.~\ref{AP0}). For other rational values of $c$ we have calculated
$F_S$ numerically. We find the following properties: (i) The maximal
static friction $F_S$ as a function of $b$ increases strictly
monotonically. (ii) Near zero it scales like 
\begin{equation}
  F_S\sim b^Q.
\end{equation}
(iii) For $b\to\infty$ it approaches $F_S^{max}$. (iv) For fixed
values of $b$ it decreases with decreasing commensurability of $c$.
More precisely: If we organize the rational numbers in a Farey 
tree\cite{remSTA.2}, we find that the static friction of some
rational value of $c$ is less than the static friction of the parents
of this rational number in the Farey tree. 

Figure~\ref{f.FS} shows $F_S(b)$ for a sequence of rational values of
$c$ which approaches the golden mean $(\sqrt{5}-1)/2$. We see that
the sequence of curves converges to a curve with a singularity at the
point of breaking of analyticity $b_c^S$. For $b<b_c^S$ the static
friction $F_S$ is exactly zero because the analyticity of the hull
function (\ref{F0.gS}) makes it possible to shift the surface atoms
of the upper body smoothly over the corrugated surface of the lower
body. Thus a Goldstone mode exists and no force is needed for
sliding. As static friction and the breaking of analyticity always
appear together, we change the meaning of $b_c^S$.  In the following
$b_c^S$ denotes the threshold above which the static friction is not
zero. In the commensurate case $b_c^S$ is equal to zero.

For $b>b_c^S$, $F_S$ increases strictly monotonically with $b$. 
Numerically we find a power law near $b_c^S$
\begin{equation}
 F_S\sim (b-b_c^S)^\alpha,\quad{\rm with}\quad\alpha\approx 2.
 \label{STA.powlaw}
\end{equation}
This behavior of $F_S$ is similar to what is found for the depinning
force of the ground state of the FK model\cite{pey.83}, except that
for the FK model $\alpha\approx 3$. We have checked that in the FKT
model the exponent $\alpha$ remains roughly two for values of
$\kappa$ down to $0.01$. Thus $\alpha(\kappa\to 0)$ is possibly not
continuous at $\kappa=0$.

\section{\protect\label{KIN}Kinetic friction in the quasistatic limit}

In this section we analyze the case of vanishing kinetic friction in
the quasistatic limit, i.e., sliding velocity $v\to 0$. The kinetic
friction for arbitrary sliding velocities will be calculated in part
II.

The kinetic friction is the lateral force which is necessary to keep
the upper body at a constant velocity $v$. Thus, the block position
$x_B$ is the control parameter contrary to Secs.~\ref{F0} and
\ref{STA} where the applied force $F$ was the control parameter. 

Quasistatic sliding gives the system enough time to relax into a
stable state after an infinitesimal change of $x_B$. The stationary
states, which are solutions of (\ref{STA.eqj}), depend on $x_B$. From
(\ref{STA.eqB}) follows the lateral force 
\begin{equation}
  F(x_B)=-\kappa\sum_{j=1}^N\xi_j(x_B)
  \label{KIN.F}
\end{equation}
which is necessary to keep the upper body at position $x_B$. In the
field of nanotribology\cite{sin.92} this lateral force is often
called friction force (for example in the term `friction force
microscopy') even though it is not of dissipative character as we
will see below.  Kinetic friction has to be defined from the {\em
work\/} which is necessary to slide two interacting surfaces over a
finite distance. Thus the kinetic friction $F_K$ is the {\em averaged
lateral force\/}, i.e.,
\begin{equation}
  F_K=-\lim_{x_B\to\infty}\frac{\kappa}{x_B}\int_0^{x_B}
      \sum_{j=1}^N\xi_j(x)\,dx.
  \label{KIN.FK}
\end{equation}

As already mentioned in the introduction, frictionless sliding has
attracted some attention in the literature. We emphasize that zero
friction for a {\em finite\/} velocity is not possible in open
systems. The kinetic friction may vanish only for sliding velocity
$v\to 0$.  Frictionless sliding simply means that Coulomb's law of
constant kinetic friction is replaced by some other (e.g., viscous)
friction law where the kinetic friction vanishes for $v\to 0$. 

What is the condition for vanishing kinetic friction in the
quasistatic limit? The answer is: The state of the system must follow
the quasistatic change of $x_B$ adiabatically. In this case the
$\xi_j$'s are functions of $x_B$. Thus the potential energy
(\ref{MOD.V}) with $F\equiv 0$ becomes a continuous periodic function
of $x_B$, i.e., 
\begin{displaymath}
	\tilde{V}(x_B)\equiv V(\xi_1(x_B),\ldots,\xi_N(x_B),x_B).
\end{displaymath}
Using (\ref{STA.eq}) we get for the lateral force (\ref{KIN.F})
\begin{displaymath}
	F(x_B)=\partial_{x_B}V=\frac{d\tilde V}{dx_B},
\end{displaymath}
which is a conservative force oscillating around zero. Its
average is zero, i.e., $F_K=0$. 

As long as $b$ is small there exists only one solution of
(\ref{STA.eqj}) for any value of $x_B$. It is given by the same hull
function $g_S$ introduced in Sec.~\ref{F01} in order to characterize
the ground state (\ref{F0.gS}) of the undriven system even in the
commensurate case. The only difference is that now $x_B$ is
arbitrary.  For $b>b_c^K$ additional stationary states appear as
already noted in Sec.~\ref{F02}. The reason is that the hull function
$g_S$ is no longer uniquely defined because the solutions of
(\ref{F0.gS.eq}) develop loops at values of $x$ which are integer
multiples of $1/Q$. The largest loop appears at $x=0$ (see also
Fig.~\ref{f.hull2}). The loops correspond to bistability intervals
which are bounded by saddle-node bifurcations. When we start with
some state, the quasistatic change of $x_B$ moves the state
adiabatically until a saddle-node bifurcation is reached. At that
point the system has to rearrange itself into a new metastable state.
During this rearrangement the particles have to move over a finite
distance where they dissipate the energy which is determined by the
difference between the potential energy before and after the 
rearrangement. Thus $F_K\not=0$ for $b>b_c^K$. As additional states
and kinetic friction in the quasistatic limit always occur together,
$b_c^K$ also denotes the threshold below which the kinetic friction
vanishes for $v\to 0$.

For $c=1$ and $1/2$ we are able to calculate $b_c^K$ analytically 
(see App.~\ref{AP2}). For other values of $c$ this is possible only
numerically. Figure~\ref{f.bck} shows numerical results of $b_c^K$ 
as a function of $\kappa$ for a sequence of rational values of $c$ 
approaching the golden mean. We find the following properties: (i) 
$b_c^K$ as a function of $\kappa$ increases strictly monotonic. 
(ii) Near zero it scales like
\begin{equation}
  b_c^K\sim \kappa^{1/Q}.
\end{equation}
(iii) In the Tomlinson limit $\kappa\to\infty$ it approaches 
$\kappa/(2\pi)$. (iv) For fixed values of $\kappa$ it increases 
with decreasing the commensurability of $c$. 

In the incommensurate limit $b_c^K$ becomes finite in the 
limit $\kappa\to 0$. The value $b_c^K(\kappa\to 0)$ is $\approx 0.154$ 
which can not be distinguished numerically from the value of the
point of breaking of analyticity of the FK model for the golden 
mean\cite{pey.83}. Thus $b_c^K(\kappa\to 0)=b_c^K(\kappa=0)$. 
Numerically we find that $b_c^K$ scales near $\kappa=0$ like
\begin{equation}
 b_c^K-b_c^K(0)\sim \kappa^\beta,\quad{\rm with}\quad
   \beta\approx\frac{1}{2}.
 \label{KIN.bck}
\end{equation}

In the incommensurate case the breaking of analyticity of the hull
function corresponds to the appearance of loops in the solution of
(\ref{F0.gS.eq}).  Thus, $b_c^K=b_c^S$ which is confirmed by our
numerical results.  Therefore both static friction and kinetic
friction in the quasistatic limit are zero. In the commensurate case,
on the other hand, zero kinetic friction does not imply zero static
friction. In fact $F_S$ is always unequal to zero.

\section{\protect\label{CON}Conclusion}

In this paper we have introduced a simple model of wearless friction,
the Frenkel-Kontorova-Tomlinson (FKT) model (see Fig.~\ref{f1}),
a combination of the Frenkel-Kontorova (FK) model and the 
Tomlinson model. Here we have investigated only the static properties
of the FKT model. The behavior of the system strongly depends on (i) 
the strength $b$ of the interaction between the sliding surfaces and 
(ii) the commensurability of the surface lattices. There are three 
threshold values of $b$ at which the behavior changes qualitatively. 

The first threshold is denoted by $b_c^S$. Below $b_c^S$ the static
friction is zero whereas above $b_c^S$ it increases like a power law
and approaches $F_S^{max}=bN$ in the asymptotic limit. In the
commensurate case where the ratio $c$ of the lattice constants of
both sliding surfaces is rational $b_c^S$, is zero. The exponent of
the power law is given by the denominator of $c$. In the
incommensurate case, $b_c^S$ is finite and increases with the ratio
$\kappa$ of the stiffnesses of the leaf and coil springs (see 
Fig.~\ref{f1}). For the golden mean the exponent is roughly two.

For $b$ below the second threshold denoted by $b_c^K$, the kinetic
friction $F_K$ is zero in the limit of quasistatic sliding (i.e.,
sliding velocity $v\to 0$). That is, for $b<b_c^K$ the kinetic
friction behaves similar to viscous friction. For $b>b_c^K$
Tomlinson's basic dissipation mechanism leads to a finite kinetic
friction. Therefore $b_c^K$ is also the threshold above which
additional stationary states appear which are stationary for fixed
relative positions of the sliding bodies.  The threshold $b_c^K$
increases with $\kappa$ and is always larger than zero. Therefore, in
the commensurate case vanishing kinetic friction does not imply
vanishing static friction.  The FKT model for $b>b_c^K$ is an example
of a dry-friction system which dynamically behaves like a viscous
fluid under shear even though the static friction is not zero. This
is a generic behavior which strongly violates Coulomb's third law.

The third threshold denoted by $b_c^m$ is important for the precise 
meaning of the maximal static friction $F_S$. Below this threshold 
the ground state of the undriven FKT model is the only stable state.
For $b>b_c^m$ additional metastable states appear. The first
metastable state is characterized by equally sized domains. The
domain states are almost the ground state but for different virtual
positions of the sliding bodies (for an example, see
Fig.~\ref{f.c1kink}). The bulk domain states differ from the exact
ground state because the real position is the average of the virtual
ones. In the commensurate case the gap in the energy between the
first metastable state and the ground state is an {\em extensive\/}
quantity, i.e., it is proportional to the number of particles. This
is a consequence of the {\em nonlocal\/} character of the FKT model. 

Since the static friction is the force which is necessary to start
sliding, it depends on the state of the system. Therefore static
friction is not uniquely defined for $b>b_c^m$, but the function
$F_S(b)$ introduced in Sec.~\ref{STA} gives an upper limit of the
actual static friction. It is defined as the smallest force above
which there exist no stable state. We have seen that $F_S$ is given
by the force required to depin the model from the ground state.

The relation between these thresholds depend on whether $c$ is
rational or irrational. For the commensurate case the relation is
\begin{equation}
  0=b_c^S<b_c^K<b_c^m.
  \label{CON.bc.c}
\end{equation}
For the incommensurate case it is
\begin{equation}
  0<b_c^S=b_c^K=b_c^m.
  \label{CON.bc.ic}
\end{equation}

Concerning the static properties of the FKT model there remain a lot
of open questions to be investigated. From the mathematical point of
view, a number of conjectures has been presented in this paper for
which we have strong numerical evidence but no rigorous proof. For
example, it is evident from physical intuition that the ground state
has the largest depinning force. This has been found also in more
complex system\cite{rob.95}. It would be desirable to have a proof of
this conjecture at least for a certain class of models.

Although the FKT model is a very simplified model even for atomically
flat surfaces, we believe that it is able to mimic some of the
qualitative behavior of real systems. The FKT model has the potential
to become a prototypical model for dry friction of atomically flat
surfaces like the FK model for the commensurate-incommensurate
transition and charge-density waves.

The FKT model can be extended in various directions. An important one
would be the introduction of quenched randomness due to unavoidable
surface impurities because in FK-like models it may destroy the 
frictionless state in the incommensurate case\cite{mat.94}.

\section*{Acknowledgment}

We gratefully acknowledge helpful discussions with T. Baumberger, T.
Gyalog, M.~O. Robbins, R. Schilling, T. Strunz, and H. Thomas. 
We especially thank H. Thomas for careful reading of the manuscript. 
This work was supported by the Swiss National Science Foundation.

\appendix

\section{\label{AP1}The first metastable state in the Tomlinson limit
    for \protect\lowercase{{$c=1$}}}

For the Tomlinson limit (i.e., $\kappa,b\to\infty$ but $b/\kappa$
finite) of the undriven FKT model we drop the first two terms of the
potential (\ref{MOD.V}). Stationary states are solutions of
\begin{equation}
  \partial_{x_j}V=\kappa\xi_j-b\sin 2\pi(x_B+cj+\xi_j)=0
  \label{AP1.xi}
\end{equation}
and (\ref{F0.xisum}).
A stationary states is stable if the second derivative of the
potential $DV$ is a positive definite matrix. It reads
\begin{equation}
 DV\equiv\left(\begin{array}{ccccc}D_1&0&\ldots&0&-\kappa\\
   0&D_2&\ldots&0&-\kappa\\\vdots&\vdots&\ddots&\vdots&\vdots\\0&0&
   \ldots&D_N&-\kappa\\kappa&-\kappa&\ldots&-\kappa&N\kappa
   \end{array}\right),
  \label{AP1.DV}
\end{equation}
where
\begin{equation}
  D_j=\kappa+C_j,\quad{\rm with}\quad C_j=-2\pi b\cos 
   2\pi(x_B+cj+\xi_j).
  \label{AP1.D}
\end{equation}
It is positive definite if all $D_j$'s and $\mbox{det}\,DV$ are
positive. After some algebraic operations one finds
\begin{equation}
  \mbox{det}\,DV=\kappa\left(\prod_{j=1}^ND_j\right)\sum_{j=1}^N
   \frac{C_j}{D_j}.
  \label{AP1.det}
\end{equation}
Thus the stability conditions read
\begin{equation}
  \sum_{j=1}^N\frac{C_j}{D_j}>0\quad\mbox{and}\quad D_j>0,
   \ j=1,\ldots,N.
  \label{AP1.stab}
\end{equation}
A sufficient condition for stability is
\begin{equation}
  C_j>0,\ j=1,\ldots,N.
  \label{AP1.suf.stab}
\end{equation}

For the case $c=1$ we investigate the existence and stability of
solutions containing a kink and an antikink in more detail. The
solution has two domains where $\xi$ is uniform in each
domain. There are $N_\pm$ particles with $\xi_\pm$ fulfilling
(\ref{AP1.xi}) with $N_++N_-=N$. The definition $p=N_+/N$ measures
the extent of the plus domain relative to the total extent of the
system. Because of (\ref{F0.xisum}) it has to fulfill
\begin{equation}
  p\xi_++(1-p)\xi_-=0.
  \label{AP1.p}
\end{equation}
The solution is stable if $pC_+/D_++(1-p)C_-/D_->0$ and $D_\pm>0$.
Eqs.~(\ref{AP1.xi}), (\ref{AP1.p}), and
\begin{equation}
  p\frac{C_+}{D_+}+(1-p)\frac{C_-}{D_-}=0
  \label{AP1.sum}
\end{equation}
define implicitly $b/\kappa$ as an even function of $p-1/2$. 
The numerical solution is shown in Fig.~\ref{f.c1kink}. The
minimum at $p=1/2$ can be calculated analytically because
(\ref{AP1.p}) implies $\xi_-=-\xi_+$. Furthermore $x_B=0$ because of
(\ref{AP1.xi}) and therefore $C_\pm=-2\pi b\cos 2\pi\xi_+$. Thus
(\ref{AP1.sum}) reduces to $C_\pm=0$ which implies $\xi_+=1/4$. With
(\ref{AP1.xi}) we get $b_c^m=\kappa/4$.

\section{\protect\label{AP0}$F_S$ for \lowercase{$c=1$} and 
\lowercase{$c=1/2$}}

In accordance with our conjecture that the static friction $F_S$ is 
given by the depinning force of the ground state, we solve the
delay equation (\ref{F0.gS.eq}) for the hull function $g_S$. Using
(\ref{F0.gS}) and (\ref{STA.eqB}) we get
\begin{equation}
  F(x_B)=-\kappa\sum_{j=1}^N g_S(x_B+cj).
  \label{AP0.F}
\end{equation}
The static friction $F_S$ is determined by the maximum of $F(x_B)$ 
which can be obtained by equating the extrema.

First we calculate $F_S$ for $c=1$. Because of the periodicity of 
$g_S$ the delay equation turns into an algebraic equation
\begin{equation}
  0=-\kappa g_S(x)+b\sin 2\pi[x+g_S(x)].
  \label{AP0.eq1}
\end{equation}
For $F$ we get
\begin{equation}
  F(x_B)=-\kappa N g_S(x_B)=-bN\sin 2\pi[x+g_S(x)].
  \label{AP0.F1}
\end{equation}
The extrema of $F$ are also extrema of $g_S$. From (\ref{AP0.eq1}) 
follows
\begin{equation}
  \frac{dg_S}{dx}=\frac{2\pi b\cos 2\pi[x+g_S(x)]}{\kappa-2\pi b\cos 
	2\pi[x+g_S(x)]}.
  \label{AP0.gs1}
\end{equation}
The extrema are implicitly defined by $x+g_S=n/2+1/4$, $n$ integer, 
which leads to
\begin{equation}
  F_S=bN.
  \label{AP0.FS1}
\end{equation}
Thus $F_S$ is equal to the overall upper bound $F_S^{max}$
which proves our conjecture that the ground state has the largest
depinning force.

For the case $c=1/2$ we introduce $g_1(x)\equiv g_S(x)$ and $g_2(x)
\equiv g_S(x+1/2)$. In accordance with the symmetries
(\ref{F0.gS.sym}) the delay equation (\ref{F0.gS.eq}) turns into a
system of two coupled algebraic equations 
\begin{equation}\begin{array}{rcl}
  0&=&2g_2(x)-(2+\kappa)g_1(x)+b\sin 2\pi[x+g_1(x)]\\
  0&=&2g_1(x)-(2+\kappa)g_2(x)-b\sin 2\pi[x+g_2(x)].
\end{array}
\label{AP0.eq2}
\end{equation}
For $F$ we get
\begin{equation}
  F(x_B)=-\frac{\kappa N}{2}[g_1(x_B)+g_2(x_B)].
  \label{AP0.F2}
\end{equation}
It is not possible to give an analytic expression of $F_S$ as an 
explicit function of $b$. But we are able to express $F_S(b)$ in 
parametric form, i.e., $F_S(D)$ and $b(D)$ where the parameter $D$ is
defined by
\begin{equation}
  D=g_1-g_2.
  \label{AP0.D}
\end{equation}
Together with (\ref{AP0.F2}) we can rewrite (\ref{AP0.eq2}):
\begin{equation}\begin{array}{rcl}
  0&=&\frac{F}{N}+b\cos 2\pi\left(x-\frac{F}{\kappa N}\right)\cdot
	\sin \pi D\\
  0&=&-\left(2+\frac{\kappa}{2}\right)\,D+b\sin 2\pi\left(x-\frac{F}
	{\kappa N}\right)\cdot\cos \pi D.
\end{array}
\label{AP0.eq2v2}
\end{equation}
Eliminating $x$ we get
\begin{equation}
  \frac{F}{F_S^{max}}=\pm\sin\pi D\sqrt{1-\left(\frac{(4+\kappa)D}
	{2b\cos\pi D}\right)^2}.
  \label{AP0.FS}
\end{equation}
In order to get the extrema of $F(x_B)$ we differentiate
(\ref{AP0.eq2v2}) in $x$. Using $dF/dx\equiv 0$, eliminating $dD/dx$
and $x-F/(\kappa N)$ we get
\begin{equation}
  b=\frac{4+\kappa}{2\pi\cos^2\pi D}\sqrt{(\pi D+\sin\pi D
   \cos\pi D)\pi D}.
  \label{AP0.b}
\end{equation}
Eqs.~(\ref{AP0.b}) and (\ref{AP0.FS}) define $F_S(b)$ parametrically.
Expansions into Taylor series around $D=0$ and $D=1/2$ yield
\begin{equation}
  \frac{F}{F_S^{max}}=\frac{\pi}{4+\kappa}\,b+{\cal O}(b^2),
  \label{AP0.FS.b0}
\end{equation}
and
\begin{equation}
  \frac{F}{F_S^{max}}=1-\frac{4+\kappa}{4}\,b^{-1}+{\cal O}(b^{-2}),
  \label{AP0.FS.bi}
\end{equation}
respectively.

\section{\protect\label{AP2}\protect\lowercase{$b_c^{\protect
   \uppercase{K}}$}for \protect\lowercase{$c=1$} and \protect
   \lowercase{$c=1/2$}}

Loops appear in the solutions of (\ref{F0.gS.eq}) when the slope of
$g_S(x)$ becomes infinite. For $c=1$ the delay equation
(\ref{F0.gS.eq}) turns into the algebraic equation (\ref{AP0.eq1}). 
Near $x=0$ we make the ansatz
\begin{equation}
  g_S(x)=g'x+{\cal O}(x^2)
  \label{AP2.ansatz}
\end{equation}
which yields
\begin{equation}
  0=[-\kappa g'+2\pi b(1+g')]x+{\cal O}(x^2).
\end{equation}
Equating the coefficient of $x^1$ equal to zero we get
\begin{equation}
  g'=\frac{2\pi b}{\kappa-2\pi b}.
  \label{AP2.gs}
\end{equation}
The slope of $g_S(x)$ becomes infinite at $x=0$ for $b$ equal to
\begin{equation}
  b_c^K=\frac{\kappa}{2\pi}.
  \label{AP2.bck1}
\end{equation}

For the case $c=1/2$ the delay equation (\ref{F0.gS.eq}) turns 
into the system of algebraic equations (\ref{AP0.eq2}) for
$g_1(x)\equiv g_S(x)$ and $g_2(x)\equiv g_S(x+1/2)$.  Again we expand
$g_1$ and $g_2$ in Taylor series:
\begin{equation}
  g_i=g_i'x+{\cal O}(x^2),\quad i=1,2.
  \label{AP2.ansatz2}
\end{equation}
Inserting this ansatz into (\ref{AP0.eq2}) and equating the
coefficient of $x^1$ equal to zero leads to
\begin{equation}\begin{array}{rcl}
  (2+\kappa-2\pi b)g_1'-2g_2'&=&2\pi b\\
  -2g_1'+(2+\kappa+2\pi b)g_2'&=&-2\pi b.
\end{array}
\label{AP2.eq3}
\end{equation}
The solution is
\begin{equation}
  g_{1,2}=\frac{2\pi b(2\pi b\pm\kappa)}{(4+\kappa)\kappa-(2\pi b)^2}.
  \label{AP2.g12}
\end{equation}
Thus the slope of $g_S$ becomes infinite at $x=0$ and $x=1/2$ for 
$b$ equal to
\begin{equation}
  b_c^K=\frac{\sqrt{(4+\kappa)\kappa}}{2\pi}.
  \label{AP2.bck2}
\end{equation}

\begin{references}
\bibitem{bow.54}
F.P. Bowden and D. Tabor, {\it Friction and Lubrication} (Oxford
University Press, 1954).

\bibitem{sin.92}
I.L. Singer and H.M. Pollock (eds), {\it Fundamentals of Friction:
Macroscopic and Microscopic Processes} (Kluwer Academic Publishers,
Dordrecht, 1992).

\bibitem{tom.29}
G.A. Tomlinson, Phil. Mag. Series 7, {\bf 7}, 905 (1929).

\bibitem{mcc.89}
G.M. McClelland, in {\it Adhesion and Friction}, M. Grunze and H.J.
Kreuzer (eds.), Springer Series in Surface Science {\bf 17}, 1
(Springer Verlag, Berlin, 1990).

\bibitem{mcc.92}
G.M. McClelland and J.N. Glosli, in Ref.~\onlinecite{sin.92}, p. 405.

\bibitem{fre.38}
Y.I. Frenkel and T. Kontorova, Zh. Eksp. Teor. Fiz. {\bf 8}, 1340 
(1938).

\bibitem{sok.84}
J.B. Sokoloff, Surf. Sci. {\bf 144}, 267 (1984).

\bibitem{aub.78} 
  S. Aubry, in  {\it Solitons and Condensed Matter Physics}, edited by
  A.R. Bishop and T. Schneider (Springer-Verlag, New York, 1978).

\bibitem{shi.93}
K. Shinjo and M. Hirano, Surf. Sci. {\bf 283}, 473 (1993).

\bibitem{hir.93}
M. Hirano and K. Shinjo, Wear {\bf 168}, 121 (1993).

\bibitem{bur.67}
R. Burridge and L. Knopoff, Bull. Seismol. Soc. Am. {\bf 57}, 341 
(1967).

\bibitem{car.89}
J.M. Carlson and J.S. Langer, Phys. Rev. Lett. {\bf 62}, 2632 (1989), 
Phys. Rev. A {\bf 40}, 6470 (1989).

\bibitem{sor.92}
D. Sornette, J. Phys. I (France) {\bf 2}, 2089 (1992).

\bibitem{sch.93}
J. Schmittbuhl, J.-P. Vilotte, and S. Roux, Europhys. Lett. {\bf 21},
375 (1993).

\bibitem{hes.94}
F. Heslot, T. Baumberger, B. Perrin, B. Caroli, and C. Caroli, Phys.
Rev. E {\bf 49}, 4973 (1994).

\bibitem{val.88}
F. Vallet, R. Schilling, and S. Aubry, J. Phys. C {\bf 21}, 67 (1988).

\bibitem{rab.65}
E. Rabinowicz, {\it Friction and Wear of Materials} (Wiley \&
Sons, New York, 1965).

\bibitem{remSTA.2}
The Farey tree is generated by successive subdivision of intervals.
Each interval defined by the rationals $M_L/N_L$ and $M_R/N_R$ is
divided  at $(M_L+M_R)/(N_L+N_R)$ into two subintervals. This process
starts with $[0/1,1/1]$ and is repeated to infinity.

\bibitem{pey.83}
M. Peyrard and S. Aubry, J. Phys. C {\bf 16}, 1593 (1983).

\bibitem{rob.95}
M.O. Robbins, private communication.

\bibitem{mat.94}
H. Matsukawa and H. Fukuyama, Phys. Rev. B {\bf 49}, 17286 (1994).

\end{references}

\begin{figure}
%\begin{center}
%\epsfig{file=FKT-Modell_fje.eps, width=\textwidth} 
%\end{center}
\caption{\protect\label{f1}
The Frenkel-Kontorova-Tomlinson (FKT) model where $x_B$ is the
position of the upper sliding body, $F$ is the applied force, $c$ is
the lattice constant of the surface of the upper sliding body,
$\kappa_1$ and $\kappa_2$ are the stiffness constants of the coil and
the leaf springs, respectively, $\xi_j$ is the position of particle
$j$ relative to the support of its leaf spring, and $b$ is the
strength of the interaction with the surface of the fixed rigid lower
body. Its lattice constant defines the length unit.  The stiffness of
the coil springs defines the unit of the strength of interaction.}
\end{figure}

\begin{figure}[htbp]
%\begin{center}
%\epsfig{file=hull_fje.eps, width=\textwidth} 
%\end{center}
\caption[hull]{\protect\label{f.hull}
The hull function $g_S(x)$ for $\kappa=1$ and
$c=(\sqrt{5}-1)/2\approx 144/233$. Solid line: $b=0.4$. Dotted line:
$b=0.5$.}
\end{figure}

\begin{figure}[htbp]
%\begin{center}
%\epsfig{file=hull2.eps, width=\textwidth} 
%\end{center}
\caption[hull2]{\protect\label{f.hull2}
The hull function $g_S(x)$ in the Tomlinson limit (i.e.,
$\kappa,b\to\infty$) defined by (\ref{F0.g}) with $b/\kappa=1$. The
solid line indicates stable states of a single Tomlinson oscillator
with fixed support $x$. Dashed-dotted and dotted lines indicate
metastable states and unstable states, respectively. The squares 
denote the ground state of the FKT model for $\kappa=10$, $b=10$, and
$c=21/34$.}
\end{figure}

\begin{figure}[htbp]
%\begin{center}
%\epsfig{file=fkink.eps, width=\textwidth} 
%\end{center}
\caption[kink]{\protect\label{f.kink}
The first metastable states are kink-antikink pairs. They separate
domains with different ground states which are characterized by
``virtual'' positions $x_B$ of the upper body. The upper part of the
figure gives an example for $c=1$ in the Tomlinson limit.}
\end{figure}

\begin{figure}[htbp]
%\begin{center}
%\epsfig{file=phase.eps, width=\textwidth} 
%\end{center}
\caption[c1kink]{\protect\label{f.c1kink}
Phase diagram of the elementary kink-antikink pair for
$c=1$ and $N\to\infty$. The relative portion of one domain is denoted
by $p$.}
\end{figure}

\begin{figure}[htbp]
%\begin{center}
%\epsfig{file=ffs.eps, width=\textwidth} 
%\end{center}
\caption[FS]{\protect\label{f.FS}
Static friction $F_S$ as a function of $b$ for $\kappa=1$ and for a
sequence of rational values of $c$ (thin lines) which converges to
the golden mean $(\sqrt{5}-1)/2$ (thick line).}
\end{figure}

\begin{figure}[htbp]
%\begin{center}
%\epsfig{file=b_c_fje.eps, width=\textwidth} 
%\end{center}
\caption[bck]{\protect\label{f.bck}
The threshold $b_c^K$ as a function of $\kappa$
for a sequence of rational values of $c$ (thin lines) which converges 
to the golden mean $(\sqrt{5}-1)/2$ (thick line).}
\end{figure}

\end{document}
